\section{\texorpdfstring{Převody do desítkové soustavy}{Prevody do desitkove soustavy}}

% =====================================================================================================================
\begin{frame}
    \frametitle{Převod do desítkové soustavy}
    
		\begin{minipage}[t][0.4\textheight][t]{\textwidth}
      Převodem do desítkové soustavy je jen vyjádřením hodnoty polynomu: \\
    
      \begin{itemize}
        \item \textbf{součtem vah násobených příslušnou číslicí} nebo
        \item \textbf{pomocí Hornerova schéma.}
      \end{itemize}
		\end{minipage}
		
		\begin{minipage}[t][0.5\textheight][t]{\textwidth}
      Suma součinů vah byla ukázána již na předchozích slajdech:

      \small
			\boldmath
			\begin{align*}
      (0110,0100)_2 =& 1 \cdot 2^2 + 1 \cdot 2^1 + 1 \cdot 2^{-2} \\
				=& (6,25)_{10}
			\end{align*}
      
      \normalsize
      Řád soustavy se zapisuje jako index k závorce.
		\end{minipage}
		
	\end{frame}
  
  % =====================================================================================================================
  \begin{frame}
    \frametitle{Hornerovo schéma}
    
		\begin{minipage}[t][0.4\textheight][t]{\textwidth}
      \begin{itemize}
        \item známe z algebry,
        \item rychlý převod, bez nutnosti výpočtu mocnin základu,
        \item celočíselná a neceločíselná část se převádí odděleně (dvě schéma, liší se základ, viz dále),
        \item velmi vhodné pro dvojkovou soustavu - jednoduché násobení.
      \end{itemize}
		\end{minipage}
		
		\begin{minipage}[t][0.5\textheight][t]{\textwidth}
      \boldmath
      Převádíme $(1101100)_2=(108)_{10}$
      
      \begin{center}
        \begin{tabular}{p{5mm} p{10mm} p{10mm} p{10mm} p{10mm} p{10mm} p{10mm} p{10mm} p{5mm}}
         $($ & $1$ & $1$ & $0$ & $1$ & $1$ & $0$ & $0$ & $)_2$ \\
             & $0$ & \textcolor{red}{$2\cdot$}$1$ & \textcolor{red}{$2\cdot$}$3$ & \textcolor{red}{$2\cdot$}$6$ & \textcolor{red}{$2\cdot$}$13$ & \textcolor{red}{$2\cdot$}$27$ & \textcolor{red}{$2\cdot$}$54$ & \\
         \hline \\
         $($ & $1$ & $3$ & $6$ & $13$ & $27$ & $54$ & $108$ & $)_{10}$
        \end{tabular}
      \end{center}
		\end{minipage}
		
	\end{frame}
  
  % =====================================================================================================================
  \begin{frame}
    \frametitle{Hornerovo schéma - neceločíselná část}
    Hornerovo schéma lze použít pro neceločíselnou část, ale
    
		\begin{minipage}[t][0.4\textheight][t]{\textwidth}
      \begin{itemize}
        \item násobící činitel musí být převrácenou hodnotou základu,
        \item je nutné začít odzadu (číslo je ve schéma zrcadlově),
        \item musíme počítat až do číslice za desetinnou čárkou (vždy 0), počet číslic = počet násobení,
        \item násobení necelých čísel je obtížnější.
      \end{itemize}
		\end{minipage}
		
		\begin{minipage}[t][0.5\textheight][t]{\textwidth}
      \boldmath
      Převádíme $(0,0101)_2=(0,3125)_{10}$
      
      \begin{center}
        \begin{tabular}{p{5mm} p{15mm} p{15mm} p{15mm} p{15mm} p{15mm} p{5mm}}
         $($ & $1$ & $0$ & $1$ & $0$ & $,0$ & $)_2$ \\
             & 0 & \textcolor{red}{$0{,}5\cdot$}$1$ & \textcolor{red}{$0{,}5\cdot$}$0{,}5$ & \textcolor{red}{$0{,}5\cdot$}$1{,}25$ & \textcolor{red}{$0{,}5\cdot$}$0{,}625$ & \\
         \hline \\
         $($ & $1$ & $0,5$ & $1,25$ & $0,625$ & $0,3125$ & $)_{10}$
        \end{tabular}
      \end{center}
		\end{minipage}
		
	\end{frame}